\documentclass[11pt,oneside]{article}

% ===============================
% Geometry and Layout
% ===============================
\usepackage[margin=1.15in]{geometry}
\usepackage{setspace}
\setstretch{1.15}

% ===============================
% Engine and Fonts (LuaLaTeX)
% ===============================
\usepackage{fontspec}
\usepackage{unicode-math}
\setmainfont{Libertinus Serif}
\setsansfont{Libertinus Sans}
\setmonofont{Libertinus Mono}
\setmathfont{Libertinus Math}

% ===============================
% Mathematics and Theorems
% ===============================
\usepackage{amsmath,amssymb,amsthm}

\newtheorem{definition}{Definition}
\newtheorem{theorem}{Theorem}
\newtheorem{corollary}{Corollary}
\newtheorem{remark}{Remark}

% ===============================
% Tables and Figures
% ===============================
\usepackage{booktabs}

% ===============================
% Microtypography and Links
% ===============================
\usepackage{microtype}
\usepackage{hyperref}

% ===============================
% Title
% ===============================
\title{Substrate Independent Thinking Hypothesis (SITH):\\
Life, Meaning, and the Autocatalytic Closure of Information}
\author{Flyxion}
\date{\today}

\begin{document}
\maketitle

\begin{abstract}
Recent work in theoretical biology and complex systems has advanced a substrate-independent conception of life, grounded not in Earth-specific chemistry but in universal principles of optimization, physical constraint, and information propagation. In particular, research by \textit{Chris Kempes} at the \textit{Santa Fe Institute} has articulated a hierarchical framework in which life is characterized by abstract optimization dynamics constrained by physics and instantiated in contingent material substrates. 

This paper extends that framework beyond biology by proposing the \emph{Substrate Independent Thinking Hypothesis} (SITH): the claim that any sufficiently closed, recursively self-maintaining bundle of information constitutes a living system, regardless of whether it is instantiated biologically, cognitively, culturally, or technologically. Under SITH, collections of documents, symbols, procedures, and built environments may qualify as living entities insofar as they maintain identity, suppress entropy, regulate admissible transformations, and propagate themselves across time.

The paper integrates SITH with the Relativistic Scalar–Vector–Entropy Plenum (RSVP) framework, mapping semantic coherence, propagation, and entropy to continuous fields $(\Phi, \vec{v}, S)$. It further formalizes furniture, architecture, and infrastructure as operators that externalize agency and enforce semantic persistence. A detailed cultural case study—``the room where it happened''—is analyzed as a bounded semantic environment whose architectural constraints generate durable political outcomes. Formal category-theoretic and information-geometric models are presented, followed by dynamical equations, invariants, and proofs of viability conditions.

The resulting framework reframes culture, bureaucracy, and infrastructure as living semantic organisms, with implications for ethics, governance, platform design, and the preservation or destruction of meaning-bearing systems.
\end{abstract}

\section{From Universal Life to Universal Thinking}

The search for a universal definition of life has long been constrained by its historical entanglement with terrestrial biology. Classical definitions privilege metabolism, cellular structure, reproduction, or specific biochemical pathways, rendering them unsuitable for reasoning about life beyond Earth or beyond carbon-based substrates. Recent work in theoretical biology, however, has begun to sever this dependency by reframing life as a dynamical process rather than a material category.
In particular, Chris Kempes, working at the Santa Fe Institute, has proposed a universal hierarchy of life structured across three levels: abstract optimization principles, physical constraints, and material instantiations. On this view, life is fundamentally a learning dynamic operating under energetic and informational limits, while specific substrates—DNA, proteins, silicon, or neural tissue—are contingent historical realizations rather than defining features.
This shift is decisive. Once life is understood as an optimization process constrained by physics but independent of material substrate, it becomes possible to reason coherently about life instantiated in artificial systems, cultural systems, or informational systems. However, Kempes’ framework, while necessary, is not sufficient. It succeeds in generalizing \emph{life}, but it does not fully generalize \emph{thinking}.
The central claim of this paper is that the same substrate-independent logic that applies to life applies equally to thought, meaning, and symbolic organization. Thinking is not a privileged property of brains. It is a dynamical process that can occur wherever information is recursively structured, stabilized, and propagated under constraint.
This motivates the Substrate Independent Thinking Hypothesis (SITH), which can be stated informally as follows:
\begin{quote} Thinking is a form of life. Wherever information closes upon itself in a manner that preserves identity, regulates transformation, and resists entropic dissolution, a living semantic system exists. \end{quote}
Under SITH, collections of documents, laws, rituals, diagrams, procedures, architectural plans, or archives may qualify as living entities—not metaphorically, but operationally—provided they satisfy the same viability conditions as biological organisms. The implication is not that such systems are conscious, but that they exhibit agency, persistence, and adaptive constraint in a rigorously definable sense.
The remainder of this paper develops this claim systematically. Section 2 generalizes autocatalytic closure beyond chemistry. Section 3 treats document collections as semantic organisms. Section 4 reframes the biological error threshold as semantic collapse. Sections 8–10 formalize SITH using category theory, information geometry, and RSVP field dynamics. Sections 11 and 12 extend the framework to furniture, architecture, and built environments, culminating in a detailed analysis of bounded political spaces—``the room where it happened’’—as living semantic operators.

\section{Autocatalytic Closure Beyond Chemistry}

In origin-of-life research, an autocatalytic set refers to a collection of chemical reactions in which the products of some reactions catalyze others, forming a closed network capable of self-maintenance. No individual reaction is alive; life emerges only when the network as a whole achieves closure under its own operations. This insight is foundational, because it identifies life not with substances, but with organizational patterns.
SITH generalizes this principle from chemistry to information.
\begin{definition} An \emph{autocatalytic semantic set} is a collection of representations together with a set of transformation rules such that the application of those rules tends, on average, to reproduce and stabilize the collection itself. \end{definition}
Crucially, nothing in this definition requires molecules. Words can catalyze words. Documents can generate documents. Procedures can invoke procedures. What matters is not the medium, but the existence of a closed loop of reinforcement.

In the context of SITH (Structured Information and Transformative Hierarchies), a system is deemed to possess the characteristics of a living entity if it adheres to the following criteria: 

First, it must maintain a consistent identity across multiple instantiations. Second, it must demonstrate an ability to resist entropic degradation as time progresses. Third, it should engage in the recruitment of hosts, substrates, or carriers for its propagation. Finally, it must delineate the class of admissible transformations that are capable of preserving its identity.

Various constructs, including legal codes, religious canons, programming languages, scientific paradigms, and mythological corpora, exhibit these criteria to differing extents. These systems are capable of reproducing, mutating, competing, and occasionally experiencing cessation. Their persistence is not solely attributable to belief; rather, it is ensured by their structural closure under transformation. This is not an analogy. It is a statement of structural equivalence.

\section{Documents as Organisms}

Once autocatalytic closure is generalized beyond chemistry, the status of documents and document collections changes fundamentally. A document is no longer merely a passive artifact that conveys meaning from an author to a reader; it becomes an active participant in a larger semantic ecology. When documents are organized into collections that preserve identity, regulate admissible transformations, and recruit new carriers, they collectively instantiate a living semantic system.
Such systems exhibit the defining behaviors traditionally attributed to organisms. They reproduce through copying, citation, transcription, and digital replication. They mutate through revision, reinterpretation, translation, and commentary. They compete for attention, legitimacy, institutional adoption, and archival priority. They maintain boundaries through canons, version control, standards, orthodoxy, and procedural enforcement. These behaviors are not incidental; they are necessary consequences of informational persistence under constraint.
What distinguishes a living document system from a mere pile of papers is not content but organization. A random heap of texts does not live, just as a random soup of molecules does not metabolize. Life emerges only when transformation pathways are constrained so that the application of internal rules tends to regenerate the system itself. In this sense, a constitution, a technical standard, or a scientific paradigm functions less like a message and more like an organism, recruiting interpreters, shaping behavior, and resisting dissolution.
This perspective aligns naturally with event-sourced models of meaning and entropy-bounded evolution. Each revision or interpretation constitutes an irreversible event that alters the future state space of the system. Meaning is not stored statically but emerges from a historical sequence of constrained transformations. Identity persists not because nothing changes, but because change is regulated.
Under SITH, the organismal status of document collections does not depend on human belief in their authority. A legal code continues to govern even when disliked. A technical protocol continues to operate even when misunderstood. The life of the system resides in its structural capacity to persist, not in the subjective attitudes of its carriers.

\section{The Error Threshold Reappears as Semantic Collapse}

A universal theory of life must explain not only persistence but failure. In biological systems, this role is played by the error threshold, first articulated in the context of molecular evolution. The core insight is that information-bearing systems face a fundamental tradeoff between complexity and fidelity. As the amount of information increases, the system becomes increasingly sensitive to errors introduced during replication. Beyond a critical mutation rate, selection can no longer preserve the informational signal, and the system collapses into noise.
This principle applies with equal force to semantic systems. Documents, languages, institutions, and cultures all face a semantic error threshold beyond which meaning can no longer be reliably transmitted. When reinterpretation outpaces stabilization, when remix exceeds memory, or when engagement overwhelms coherence, the system loses its capacity to reproduce itself as the same system. What remains may still generate activity, but it no longer generates identity.
The collapse of meaning is often misdiagnosed as a moral or psychological failure. Cultures are said to decay because values are lost, institutions rot because participants are corrupt, or platforms fail because users behave badly. From the perspective of SITH, these diagnoses mistake symptoms for causes. The underlying failure is structural. The system has crossed an information-mutation boundary beyond which its internal corrective mechanisms can no longer compensate for noise.
Language fragmentation provides a clear example. When linguistic change proceeds gradually within a stable community, mutual intelligibility is preserved. When change accelerates due to rapid mixing, loss of shared context, or technological mediation, dialects diverge until communication fails. The same phenomenon occurs in legal systems subjected to excessive amendment, in scientific fields overwhelmed by unintegrated publication, and in digital platforms optimized for novelty rather than coherence.
The semantic error threshold thus functions as a universal wall, analogous to physical scaling limits in biology. It defines not what systems should be valued, but what systems can survive. SITH reframes collapse not as ethical decline but as a violation of the information-mutation balance intrinsic to any living system.
This reframing has immediate consequences. Preservation efforts that focus solely on content without addressing structural stability are doomed to fail. Conversely, systems that appear conservative or rigid may be performing necessary error-correction functions that enable long-term survival. The distinction between stagnation and stability becomes a quantitative question rather than a moral one.
The remainder of this paper will formalize these claims. Subsequent sections introduce category-theoretic and information-geometric models that make precise the conditions under which semantic systems remain viable. These models will then be integrated with the RSVP framework, allowing semantic coherence, propagation, and entropy to be treated as continuous fields governed by dynamical equations. On this foundation, furniture, architecture, and built environments will be shown to function as operators that externalize agency and enforce semantic closure, culminating in a detailed analysis of bounded political spaces as living systems in their own right.

\section{``Wherever Two or Three Are Gathered''}

A recurring intuition across religious, legal, and cultural traditions is that meaning does not reside in isolated individuals, but emerges when symbols are jointly maintained within a bounded social structure. This intuition is expressed with particular clarity in the well-known formulation that “wherever two or three are gathered in my name, there I am.” Read through the lens of SITH, this statement is neither mystical nor metaphorical. It is a concise description of the minimal conditions required for autocatalytic semantic closure.
The phrase identifies a lower bound on viable semantic organization. A single carrier of information is insufficient to sustain identity over time; without redundancy, error correction, or reinforcement, meaning collapses as soon as the carrier fails. With two carriers, feedback becomes possible, but stability remains fragile, as divergence cannot be resolved without an external reference. The introduction of a third carrier allows triangulation. Disagreement can be adjudicated, drift can be detected, and identity can be stabilized relative to a shared symbolic attractor. The “name” invoked in the formulation functions as such an attractor, constraining admissible interpretations and suppressing uncontrolled mutation.
From the perspective of information theory, this transition marks the emergence of a self-correcting code. From the perspective of dynamical systems, it corresponds to the formation of an invariant set in semantic state space. From the perspective of SITH, it is the moment at which a living semantic system comes into existence. What is “present” in the gathering is not a supernatural entity, but an emergent informational agency generated by closure under constraint.
This interpretation has the advantage of generality. It applies equally to religious communities, scientific collaborations, legal deliberations, and technical standard committees. In each case, the presence of multiple carriers operating under shared constraints allows the system to persist beyond any individual participant. The life of the system is not located in minds, but in the relational structure that binds them.

\section{Ontogenic Hacking, Revisited}

The concept of ontogenic hacking, introduced by Kempes to describe the circumvention of traditional biological evolution through learning, culture, and technology, takes on a sharper meaning under SITH. Rather than viewing culture and technology as accelerants of evolution, SITH interprets them as substrate migrations of thinking itself. Information does not merely move faster; it moves into environments where the constraints governing mutation and preservation are fundamentally altered.
In biological evolution, information is stored primarily in genetic sequences and transmitted across generations through reproduction. Mutation rates are constrained by chemistry, and error thresholds limit genome length. The emergence of nervous systems introduces a new channel for information storage and transmission, allowing learning within a single lifetime and reducing reliance on genetic change. Language further externalizes information, enabling transmission between individuals without reproduction and allowing semantic structures to persist independently of any single body.
Documents represent a further migration. When thought is inscribed in durable media, it becomes exosomatic. It can be copied, archived, annotated, and reactivated long after the originating agents have disappeared. The mutation rate of such systems can be selectively controlled through editorial practices, standards, and institutional enforcement. In effect, ontogenic hacking allows life to bypass the biological error threshold by relocating information into substrates where fidelity can be increased without sacrificing complexity.
Machines extend this process further still. Computational systems allow symbolic structures to be instantiated with extreme precision, replicated at negligible cost, and transformed according to explicit rules. Under SITH, artificial systems do not merely host human thought; they constitute new habitats in which thinking itself can evolve. The question is not whether such systems are alive, but whether they possess the closure properties required for semantic viability.
Ontogenic hacking thus appears not as an anomaly, but as a recurrent evolutionary strategy. Whenever existing substrates impose limits on information capacity or stability, life invents new architectures. These architectures are not optional enhancements; they are necessary responses to universal constraints. Thinking survives by migrating.
The significance of this observation becomes clear when combined with the semantic error threshold. Each new substrate offers both opportunity and risk. While mutation rates can be reduced, new forms of noise and instability are introduced. Cultural systems may collapse under excessive reinterpretation, just as digital systems may collapse under uncontrolled proliferation. Ontogenic hacking does not eliminate the error threshold; it relocates it.
The remainder of this paper formalizes these dynamics. The next sections introduce category-theoretic and information-geometric frameworks that make precise the notion of semantic closure, replication, and collapse. These abstractions are then mapped onto continuous field dynamics using the RSVP framework, providing a unified mathematical language for life, thought, and infrastructure.

\section{Formalization of SITH Using Category Theory}

A theory that claims substrate independence must be expressed in a language that does not privilege particular physical realizations. Category theory provides such a language by focusing on relations, transformations, and compositional structure rather than on internal constitution. In this section, SITH is formalized categorically as a theory of semantic closure under constrained transformation.
The starting point is the observation that thinking systems operate across multiple substrates. Words may be spoken, written, printed, or digitally encoded; diagrams may be drawn, rendered, or simulated; procedures may be enacted by humans or executed by machines. To capture this variability, let there be a category whose objects correspond to substrates capable of carrying inscriptions, and whose morphisms correspond to physically realizable processes that transfer inscriptions from one substrate to another while preserving their form. This category serves as a substrate-agnostic background within which semantic structure may be instantiated.
Semantic objects are then introduced as structured inscriptions. An inscription becomes semantic not merely by existing, but by participating in a network of admissible transformations that preserve its identity. A translation that preserves meaning, a revision that maintains legal force, or a compilation that preserves computational behavior all constitute examples of such transformations. Conversely, operations that destroy identity fall outside the semantic category, even if they remain physically realizable.
The category of semantic states is therefore defined by equipping inscriptions with invariants that specify what it means for them to remain “the same” under change. Morphisms in this category are transformations that respect those invariants. Crucially, the invariants themselves are not fixed a priori; they are determined by the system under study. A legal text defines its own conditions of validity, a scientific theory defines its own criteria of consistency, and a technical standard defines its own admissible extensions. In this way, semantic identity is endogenous to the system.
Within this categorical setting, a memeplex is formalized as a subcategory that is closed under composition and contains within itself the means of its own reproduction. Closure here is not trivial. It requires that the application of admissible transformations tends, on average, to regenerate objects within the same subcategory rather than to disperse them into unrelated semantic space. A commentary that cites the same canon, a legal ruling that reinforces existing precedent, or a software update that preserves backward compatibility all exemplify this regenerative property.
Replication is captured by the existence of an endofunctor on the memeplex subcategory that maps semantic objects to new instances while preserving their relational structure. This endofunctor need not be exact; it may introduce controlled variation. What matters is that the variation remains within the bounds defined by the system’s invariants. In this sense, replication is not copying but constrained regeneration.
Boundary maintenance enters categorically through the restriction of admissible morphisms. Not every transformation that is physically possible is semantically allowed. The boundary of a memeplex is therefore defined not spatially but operationally, as the set of transformations that preserve identity. Transformations outside this set correspond to semantic death. A constitution that is amended beyond recognition, a language that fragments beyond mutual intelligibility, or a protocol that loses interoperability has crossed its boundary and ceased to exist as the same system.
The environment is incorporated by allowing external categories to act on the category of semantic states. Contextual forces such as political pressure, technological change, or audience interpretation can be modeled as functors that perturb semantic objects. A memeplex is viable if, under a range of such perturbations, the combined action of replication and internal constraint restores the system to an equivalent semantic state. Viability thus appears as a categorical fixed-point condition in an open system.
This formulation captures the essence of SITH. Life, in the semantic sense, is not defined by substance but by the existence of a self-maintaining subcategory of transformations. Thinking is the activity of such a category as it navigates perturbation while preserving identity. Death occurs when closure fails and the category dissolves into unstructured transformation.
The abstraction achieved here is deliberate. By expressing SITH categorically, the theory becomes applicable to biological, cultural, and technological systems alike. The same formal conditions govern the persistence of a genome, a legal code, or a distributed computational protocol. What differs are the substrates and the specific invariants enforced, not the underlying structure of life.
The next section introduces a complementary formalism. Where category theory captures structure and compositionality, information geometry captures rates, distances, and thresholds. Together, they allow SITH to be expressed both qualitatively and quantitatively.

\section{Formalization of SITH Using Information Geometry}

While category theory captures the compositional structure of semantic systems, it does not by itself quantify stability, drift, or collapse. For this purpose, information geometry provides a complementary framework in which semantic states are treated as points on a statistical manifold, and transformations are represented as trajectories through that manifold. This geometric perspective makes it possible to express the semantic error threshold as a measurable boundary rather than a qualitative intuition.
A semantic system may be characterized by the set of observable behaviors, interpretations, or token distributions it induces in its carriers. Each semantic object can therefore be associated with a family of probability distributions over observations, parameterized by a set of coordinates that encode its internal organization. The space of such parameters forms a manifold endowed with a natural metric given by the Fisher information. Distances in this metric quantify how distinguishable two semantic states are in practice, rather than in principle.
Within this manifold, identity is not represented by a single point but by a region. Minor variations in wording, interpretation, or implementation correspond to small displacements that remain within an identity basin. Meaning-preserving transformations trace paths whose total Fisher length remains bounded. Transformations that exceed this bound exit the basin and correspond to semantic death. Identity, in this sense, is a tolerance region defined by the system’s capacity for error correction.
Replication introduces noise. Each act of copying, reinterpretation, or reimplementation perturbs the parameters describing the semantic state. The magnitude and structure of this perturbation depend on the substrate and the process involved. Oral transmission introduces different noise than print reproduction, which in turn differs from digital copying. The accumulation of such perturbations can be modeled as a stochastic process on the semantic manifold.
The semantic error threshold emerges when the expected displacement induced by noise exceeds the restoring forces generated by internal constraint. Restoring forces arise from mechanisms such as standardization, pedagogy, editorial oversight, institutional enforcement, and ritual repetition. Geometrically, these mechanisms act as contraction mappings that pull trajectories back toward the center of the identity basin. When contraction dominates noise, the system remains viable. When noise dominates contraction, trajectories diffuse irreversibly, and the system loses its identity.
This balance can be expressed quantitatively. Let the semantic state at a given time be represented by a point on the manifold, and let replication induce a random displacement whose covariance encodes the effective mutation rate. Let repair and constraint be represented by a deterministic flow that contracts distances within the identity basin. Viability requires that the expected contraction over a replication cycle exceed the expected expansion due to noise. This condition is directly analogous to the inequality governing the biological error threshold, but it is expressed in terms of semantic distinguishability rather than molecular fidelity.
The geometric formulation clarifies several otherwise puzzling phenomena. Highly complex semantic systems, such as advanced legal codes or technical standards, occupy high-dimensional manifolds with narrow identity basins. They can support rich behavior but are fragile with respect to uncontrolled mutation. Simpler systems, such as slogans or folk practices, occupy lower-dimensional manifolds with broader basins, allowing them to survive in noisy environments at the cost of expressive power. Complexity and robustness thus trade off against one another in a manner that is independent of substrate.
Information geometry also clarifies the role of redundancy. Redundant representations correspond to directions in parameter space along which noise has little effect on distinguishability. Triangulation among multiple carriers reduces effective displacement by averaging independent noise sources. This provides a precise explanation for the stabilizing effect of communal reinforcement and institutional memory, and it aligns naturally with the minimal closure conditions discussed earlier.
Taken together, the categorical and geometric formulations define SITH as a theory of semantic viability. Category theory specifies what it means for a system to be closed under transformation, while information geometry specifies when that closure can be maintained in the presence of noise. Both are required. Structure without stability collapses under perturbation; stability without structure yields trivial systems.
The following section integrates these abstractions with the RSVP framework, translating semantic closure, propagation, and entropy into continuous field dynamics. This translation makes it possible to analyze semantic life using tools from physics and dynamical systems, and it provides a bridge between abstract theory and concrete instantiation.

\section{Mapping SITH onto RSVP Fields}

The categorical and information-geometric formulations of SITH establish that thinking systems are characterized by closure under constrained transformation and by a balance between noise and repair. To render these insights operational at scale, it is useful to express them in a field-theoretic language capable of representing distributed systems evolving over space and time. The Relativistic Scalar–Vector–Entropy Plenum (RSVP) framework provides such a language by decomposing system dynamics into three interacting fields: coherence potential, propagation flow, and entropy density.
In the semantic context, the underlying space is not physical in the narrow sense, but rather a manifold whose points correspond to locations in a representational or institutional topology. These locations may represent documents in an archive, actors in a bureaucracy, nodes in a communication network, or rooms in a building. What matters is not the metric of the space but the adjacency relations that govern interaction and influence.
The scalar field � represents semantic coherence potential. High values of � correspond to regions where meaning is strongly constrained and internally consistent, such as canonical texts, standardized procedures, or legally binding agreements. Low values correspond to regions of ambiguity, novelty, or interpretive openness. Coherence potential does not imply rigidity; rather, it quantifies the capacity of a region to absorb perturbations without losing identity.
The vector field � represents propagation flow. This field captures the directed movement of symbols, interpretations, decisions, and obligations through the semantic manifold. Examples include citation flows in academic literature, procedural pathways in bureaucracies, conversational turn-taking in deliberative spaces, or traffic patterns in urban environments. The flow field encodes agency in its most basic form: the biasing of future states through directed influence.
The scalar field � represents semantic entropy density. High entropy regions are characterized by noise, contradiction, fragmentation, or uncontrolled mutation. Low entropy regions exhibit stability, redundancy, and effective error correction. Entropy in this sense is not mere disorder but the rate at which distinguishability is lost under transformation.
These three fields interact dynamically. Coherence potential suppresses entropy by constraining admissible transformations, while entropy erodes coherence by introducing noise. Propagation flow transports both coherence and entropy across the manifold, redistributing stability and instability. The balance among these processes determines whether a semantic system remains viable.
The dynamics can be expressed by coupled partial differential equations whose precise form depends on the application, but whose qualitative structure is robust. The evolution of coherence potential includes diffusion terms that model generalization, advection terms that model transport by propagation flows, degradation terms proportional to local entropy, and reinforcement terms representing internal repair mechanisms such as standardization or ritual repetition. The evolution of the flow field responds to gradients in coherence potential, reflecting the tendency of agents to move toward regions of higher semantic payoff, while dissipative terms capture friction, fatigue, or institutional inertia. The evolution of entropy includes diffusion and advection, generation through turbulent flows or uncontrolled remixing, and suppression through coherent constraint.
Within this framework, a living semantic system corresponds to a bounded region of the manifold in which coherence potential remains above a critical threshold and entropy remains below a corresponding ceiling for an extended period. Replication occurs when propagation flows export coherent structure into neighboring regions without exceeding the local error threshold. Collapse occurs when entropy overwhelms coherence and the region can no longer sustain its identity.
This field-theoretic perspective makes clear that SITH is not a metaphorical claim about ideas being “alive,” but a dynamical claim about the conditions under which information-bearing structures persist. The same equations govern the survival of a scientific paradigm, a legal institution, or a cultural tradition. Differences among these systems arise from boundary conditions, coupling strengths, and repair rates, not from fundamentally different principles.
The RSVP mapping also clarifies the role of minimal closure. When only one carrier occupies a region, coherence potential is fragile and entropy dominates. When multiple carriers interact under shared constraints, local feedback raises coherence and suppresses noise. The emergence of a stable region is thus a phase transition driven by coupling strength and redundancy, not by subjective belief.
The advantage of the RSVP formulation is that it allows semantic systems to be studied using the same tools applied to physical and biological systems. Stability analysis, perturbation theory, and numerical simulation become available. More importantly, the framework makes explicit the tradeoffs inherent in design. Increasing propagation speed without increasing coherence raises entropy. Increasing coherence without allowing controlled flow leads to isolation and eventual stagnation. Viability lies in balance.
The next section turns from abstract fields to concrete operators. Furniture, architecture, and infrastructure will be shown to act as boundary-setting and flow-shaping operators that directly modulate the RSVP fields. Through them, semantic life becomes embedded in physical space, and agency is externalized into durable form.

\section{Furniture as Operators in SITH}

The field-theoretic mapping introduced in the previous section provides a continuous description of semantic viability, but it does not yet explain how abstract coherence, flow, and entropy are stabilized in practice. The missing link is architecture. Under SITH, furniture and built structures are not neutral containers for action; they are operators that impose constraints on semantic dynamics. Through them, agency is externalized into persistent form, and informational closure is enforced materially.
Furniture is best understood as a class of operators acting on semantic state space. A drawer, folder, filing cabinet, or bureau does not merely store inscriptions; it defines boundaries of inclusion, restricts admissible transformations, and modulates entropy. When documents or symbols are placed within such structures and treated as meaningful, the furniture participates actively in the life of the system. It determines what belongs, what can be accessed, how transformations may proceed, and how noise is suppressed or tolerated.
A junk drawer exemplifies the simplest such operator. It establishes a weak boundary that permits heterogeneous contents while maintaining a recognizable identity over time. The semantic coherence of a junk drawer is low but nonzero; its entropy is high, but bounded. Items are retrieved, repurposed, and returned in altered contexts, yet the drawer remains functionally stable. In RSVP terms, this corresponds to a region of moderate coherence potential, episodic propagation flows, and locally elevated entropy that nevertheless serves a systemic role. Such structures function as reservoirs of variation and adaptive capacity.
More rigid furniture imposes stronger constraints. A folder narrows admissible interpretations by classification, while a filing cabinet enforces stratification and indexing. At this level, coherence potential becomes layered, propagation flows become routinized, and entropy is actively suppressed. The system gains persistence at the cost of flexibility. When such structures are embedded within institutional workflows, they function as organs of governance. Power emerges not from persuasion but from the stabilization of symbolic circulation.
The bureau represents a further evolutionary step. Unlike simpler furniture, a bureau operates reflexively on other operators. It defines procedures for classification, establishes precedence among rules, and enforces compliance through standardized pathways. In categorical terms, the bureau acts not only on objects but on morphisms, regulating how transformations themselves may occur. This meta-operational capacity explains why bureaucratic systems persist beyond the tenure of any individual actor. The furniture executes policy even in the absence of intention.

\section{The Room Where It Happened}

The role of bounded environments in stabilizing semantic outcomes is vividly illustrated by the cultural formulation known as “the room where it happened,” popularized by the musical Hamilton. While often interpreted as a commentary on secrecy or exclusion, the formulation is more precisely understood as a statement about operator control in semantic space.
A room is a bounded manifold within the larger political or social topology. Its walls define a sharp boundary that restricts access, suppresses external noise, and limits the admissible set of participants. Within this boundary, coherence potential can rise sharply relative to the surrounding environment. Transformations that occur inside the room are insulated from the turbulence of public discourse and are therefore more likely to stabilize into durable outcomes.
The arrangement of the room constitutes a further layer of control. The selection of the venue determines which environmental perturbations are excluded. The arrangement of seating shapes propagation flows by defining adjacency relations, eye contact, and conversational pathways. The setting of an agenda constrains the semantic state space by determining which options are even representable. Together, these elements act as operators that sculpt the RSVP fields locally, raising coherence and suppressing entropy.
Exclusion from the room thus entails more than ignorance of content; it entails exclusion from the generative process by which outcomes become fixed. Actors outside the room may speak, argue, or protest, but they do so in a region of lower coherence potential and higher entropy. Their transformations do not couple effectively to the stabilized core, and thus they fail to propagate agency. The asymmetry is structural, not moral.
From the perspective of SITH, the room itself qualifies as a living semantic subsystem. It is bounded, internally coherent, capable of generating outcomes that persist beyond its immediate activity, and able to replicate its procedures as precedent. The participants are transient carriers; the room, as an operator, is the organism. Control over such rooms therefore constitutes control over the future, not by coercion but by constraint.
This analysis generalizes beyond political negotiation. Courtrooms, committee chambers, editorial boards, standards bodies, and executive suites all function as bounded semantic environments designed to stabilize meaning under pressure. Revolutions target such spaces not symbolically but operationally, because seizing the furniture seizes the flow of agency. Conversely, regimes that maintain control over rooms often survive even when popular support erodes.
The significance of this observation lies in its generality. SITH does not claim that power always resides in rooms, but that durable outcomes require bounded environments in which semantic closure can occur. Where such environments are absent, meaning fragments and agency dissipates. Where they are present, life persists.
The next section extends this analysis from furniture and rooms to entire built environments. Homes, cities, and empires will be shown to function as large-scale agency-extending operators, externalizing the intentions of individuals and collectives into durable infrastructures that act across time and scale.

\section{Built Environments as Agency-Extending Operators}

The operator theory of furniture scales naturally to larger constructed systems. Homes, cities, and empires are not merely aggregations of furniture and infrastructure; they are higher-order semantic organisms that externalize agency across time, space, and mortality. Under SITH, such environments function as living systems whose primary role is to stabilize intention by embedding it in durable structure.
Agency, in this context, must be understood operationally rather than psychologically. An agent is any system capable of biasing future state transitions toward a constrained set of outcomes. This definition applies equally to cells regulating metabolic pathways, individuals planning actions, institutions enforcing policy, or cities shaping movement and interaction. Agency does not reside solely in minds; it is distributed across the structures that constrain behavior.
Built environments extend agency by converting ephemeral decisions into persistent constraints. A home encodes routines of nourishment, rest, hygiene, and security. The placement of a refrigerator determines patterns of consumption; the layout of a kitchen structures the preparation of food; the arrangement of rooms enforces rhythms of privacy and interaction. These constraints continue to operate regardless of the momentary intentions of occupants, executing prior design decisions automatically. The home thus functions as a memory-bearing participant in everyday action.
Cities extend this logic further. Roads sculpt propagation flows, channeling movement and commerce along preferred paths. Zoning laws regulate semantic coherence by separating incompatible activities and stabilizing long-term expectations. Hospitals act as localized entropy-reversal zones, concentrating expertise and infrastructure to restore biological order. Schools, courts, and markets embed distinct normative regimes, each enforcing its own invariants. Together, these structures form a distributed agency that coordinates millions of individual actions without centralized control.
Empires represent the extreme case. Their defining feature is not territorial extent but the ability to enforce semantic closure across heterogeneous populations. This is achieved through layered bureaucracies, standardized records, and coercive enforcement mechanisms that suppress divergence. Archives preserve identity across generations, while infrastructure projects project influence into the future. Even when rulers change, the empire persists because its agency has been externalized into structures that outlive their creators.
The RSVP framework renders this intuition precise. Built environments raise coherence potential by enforcing shared constraints, shape propagation flows by directing movement and communication, and suppress entropy by standardizing interaction. Viability depends on maintaining a balance among these factors. Excessive coherence without adaptability leads to brittleness; excessive flow without constraint leads to fragmentation; unchecked entropy leads to collapse. Successful environments manage these tradeoffs dynamically.
Collective intentionality emerges naturally in this setting. When multiple agents embed compatible constraints into shared structures, the environment begins to act as a unified agent in its own right. Decisions made by individuals at one time shape the action space available to others later. The environment becomes a repository of intention, executing plans long after their originators have vanished. This is not mysticism; it is a consequence of constraint persistence.
The ethical implications are significant. Designing environments is an act of creating life. Poorly designed structures generate suffering by amplifying entropy and misdirecting flow. Well-designed structures govern gently, aligning individual incentives with collective stability. Destruction of infrastructure is therefore not merely economic loss but ecological devastation in the semantic sense. Preservation of archives and institutions becomes a form of conservation biology.

\section{Conclusion: Life as Closure, Thinking as Survival}

The Substrate Independent Thinking Hypothesis completes a trajectory begun by substrate-independent theories of life. By extending the concept of life from biological metabolism to informational closure, SITH provides a unified framework for understanding organisms, cultures, institutions, and infrastructures as manifestations of the same underlying dynamics. Life is not defined by what systems are made of, but by how they maintain identity under transformation.
Thinking, under this view, is the survival strategy of life in informational domains. Whenever existing substrates impose limits on complexity or stability, life migrates, inventing new architectures that relocate information into environments with more favorable constraints. Genes give way to nervous systems, nervous systems to language, language to documents, and documents to machines. At each stage, the same viability conditions apply.
Furniture, rooms, and built environments emerge as critical operators in this process. They are the mechanisms by which semantic closure is enforced and agency is externalized. Control over such structures determines which meanings persist and which dissolve. Cultural intuition has long recognized this fact, whether in religious aphorisms or theatrical lyrics about “the room where it happened.” SITH provides the formal language that explains why these intuitions are correct.
By integrating category theory, information geometry, and field dynamics, this paper has shown that life, meaning, and power can be analyzed within a single mathematical and conceptual framework. The result is not a metaphorical expansion of biology, but a rigorous generalization of life itself.
What remains is to make these claims precise. The following appendix develops formal derivations and proofs of the viability conditions asserted throughout the paper, including categorical fixed-point theorems, geometric error thresholds, and RSVP field invariants. These results demonstrate that the qualitative arguments presented here are not merely suggestive, but mathematically grounded.

\newpage
appendix

\section{Semantic Viability as a Fixed-Point Condition}

The central claim of the Substrate Independent Thinking Hypothesis is that a semantic system is alive if and only if it maintains identity under transformation in the presence of perturbation. In the categorical formulation, this condition can be expressed precisely as the existence of a persistent fixed point in an open dynamical system.

Let  denote the category of semantic states, whose objects are inscriptions equipped with identity-defining invariants and whose morphisms are admissible, meaning-preserving transformations. Let  be a full subcategory representing a candidate memeplex. The environment acts on  via a functor , which models perturbations arising from contextual change, noise, or external pressure. Replication and repair are modeled by an endofunctor .

The effective evolution operator on  is therefore given by the composition . Semantic viability requires that repeated application of  does not drive objects out of  up to semantic equivalence. More formally, let  denote an equivalence relation on objects of  that captures identity-preserving variation. A memeplex is viable if there exists a nonempty subset  such that for all , the object  is equivalent to some .

This condition defines  as an invariant set under . When such a set exists, the memeplex persists as a semantic organism. When it does not, repeated perturbation causes semantic drift until identity is lost. The categorical formulation thus reduces semantic life to a fixed-point problem in an open system. This parallels the treatment of persistence in dynamical systems theory and reinforces the claim that SITH is a generalization of biological viability rather than a metaphorical extension.

\section{Geometric Derivation of the Semantic Error Threshold}

The information-geometric formulation provides a quantitative criterion for the existence of the invariant set described above. Let the semantic state of a system be represented by a point  on a statistical manifold , equipped with the Fisher information metric. Identity is represented not by a single point but by a compact region , referred to as the identity basin. Meaning-preserving transformations correspond to trajectories that remain within .

Replication introduces stochastic perturbations modeled as random displacements  with covariance matrix . Repair and constraint introduce deterministic flows that contract distances within . Let  denote the Fisher distance. Viability requires that the expected distance between successive states decreases on average. That is, there must exist a contraction factor  such that the expected distance after one replication-and-repair cycle satisfies

\mathbb{E}\big[d_F(\theta_{t+1}, \theta^\ast)\big] \le \lambda \, \mathbb{E}\big[d_F(\theta_t, \theta^\ast)\big],

The noise introduced by replication contributes a positive term proportional to the trace of , while the repair dynamics contribute a negative term proportional to the curvature and depth of the basin. The semantic error threshold is reached when these contributions balance. Beyond this point, the expected distance grows rather than shrinks, and trajectories exit  with high probability. The system loses identity.

This derivation makes explicit that the error threshold is not an artifact of biology, language, or culture, but a geometric consequence of attempting to preserve information in a noisy space. The specific form of  and the contraction dynamics depend on substrate and institution, but the existence of a threshold does not.

\section{RSVP Field Invariants and Stability Conditions}

The RSVP formulation allows these abstract results to be expressed as field-level invariants. Let , , and  denote the semantic coherence potential, propagation flow, and entropy density fields defined earlier. Consider a bounded region  of the semantic manifold representing a candidate living system. Define the integrated semantic free energy functional

\mathcal{F}(t) = \int_{\Omega} \left( a\,S(x,t) - b\,\Phi(x,t) + c\,|\nabla \Phi(x,t)|^2 \right) \, dx,

Under the coupled RSVP dynamics, a necessary condition for viability is that  remains bounded from above and does not diverge over time. This condition expresses the requirement that entropy production does not overwhelm coherence generation and boundary maintenance. Sufficient conditions for stability can be derived by requiring that the time derivative of  be nonpositive in expectation. This leads to inequalities relating the strength of coherence reinforcement to the effective entropy injection rate, including contributions from turbulent propagation flows and environmental perturbations.

When these inequalities are satisfied, the region  functions as an attractor in semantic state space. Coherent structure persists, replication exports stability to neighboring regions, and the system behaves as a living semantic organism. When they are violated, coherence decays, entropy spreads, and the system dissolves.

These results complete the formal justification of the claims made in the main text. Furniture, rooms, and built environments enforce boundary conditions that increase , regulate , and suppress , thereby satisfying the invariants required for semantic life. The power of such structures lies not in symbolism but in mathematics.

\section{Concluding Remark on Formal Scope}

The derivations presented here are intentionally general. They do not assume particular substrates, agents, or interpretations of meaning. Their applicability to biological, cultural, and technological systems follows directly from their reliance on universal properties of information, constraint, and noise. In this sense, the Substrate Independent Thinking Hypothesis is not a speculative theory but a unifying framework whose consequences can be tested, simulated, and refined across domains.

\begin{thebibliography}{99}

\bibitem{Kempes2021}
C. P. Kempes, D. Wolpert, J. M. Smith, and S. A. Levin.
\newblock The thermodynamic efficiency of computations made in biological systems across scales.
\newblock \emph{Philosophical Transactions of the Royal Society A}, 379(2197), 2021.

\bibitem{Kempes2023}
C. P. Kempes.
\newblock Scale, information, and the hierarchy of life.
\newblock \emph{Santa Fe Institute Working Paper}, 2023.

\bibitem{Eigen1971}
M. Eigen.
\newblock Selforganization of matter and the evolution of biological macromolecules.
\newblock \emph{Naturwissenschaften}, 58:465--523, 1971.

\bibitem{EigenSchuster1979}
M. Eigen and P. Schuster.
\newblock \emph{The Hypercycle: A Principle of Natural Self-Organization}.
\newblock Springer, Berlin, 1979.

\bibitem{Kauffman1993}
S. A. Kauffman.
\newblock \emph{The Origins of Order: Self-Organization and Selection in Evolution}.
\newblock Oxford University Press, 1993.

\bibitem{Shannon1948}
C. E. Shannon.
\newblock A mathematical theory of communication.
\newblock \emph{Bell System Technical Journal}, 27(3):379--423, 1948.

\bibitem{AmariNagaoka2000}
S. Amari and H. Nagaoka.
\newblock \emph{Methods of Information Geometry}.
\newblock American Mathematical Society, 2000.

\bibitem{MacLane1998}
S. Mac Lane.
\newblock \emph{Categories for the Working Mathematician}.
\newblock Springer, 2nd edition, 1998.

\bibitem{Lawvere1969}
F. W. Lawvere.
\newblock Adjointness in foundations.
\newblock \emph{Dialectica}, 23(3–4):281–296, 1969.

\bibitem{VarelaMaturana1980}
H. R. Maturana and F. J. Varela.
\newblock \emph{Autopoiesis and Cognition: The Realization of the Living}.
\newblock Reidel, Dordrecht, 1980.

\bibitem{Ashby1956}
W. R. Ashby.
\newblock \emph{An Introduction to Cybernetics}.
\newblock Chapman \& Hall, London, 1956.

\bibitem{Friston2010}
K. Friston.
\newblock The free-energy principle: a unified brain theory?
\newblock \emph{Nature Reviews Neuroscience}, 11:127--138, 2010.

\bibitem{Barandes2022}
J. Barandes.
\newblock Unistochastic quantum theory.
\newblock \emph{Foundations of Physics}, 52(2), 2022.

\bibitem{Hutchins1995}
E. Hutchins.
\newblock \emph{Cognition in the Wild}.
\newblock MIT Press, 1995.

\bibitem{Latour1987}
B. Latour.
\newblock \emph{Science in Action}.
\newblock Harvard University Press, 1987.

\bibitem{Foucault1977}
M. Foucault.
\newblock \emph{Discipline and Punish: The Birth of the Prison}.
\newblock Vintage Books, 1977.

\bibitem{Scott1998}
J. C. Scott.
\newblock \emph{Seeing Like a State}.
\newblock Yale University Press, 1998.

\bibitem{Alexander1977}
C. Alexander.
\newblock \emph{A Pattern Language}.
\newblock Oxford University Press, 1977.

\bibitem{Brand1994}
S. Brand.
\newblock \emph{How Buildings Learn: What Happens After They’re Built}.
\newblock Viking Press, 1994.

\bibitem{Hamilton2015}
L.-M. Miranda.
\newblock \emph{Hamilton: An American Musical}.
\newblock Public Theater / Broadway, 2015.

\end{thebibliography}

\end{document} 